% !Mode:: "TeX:UTF-8"
% !TEX root = ../../main.tex
% ============================================================
% TOPIC 4
% ============================================================
\mashgulot{4}{Computing the rank of a matrix. Singular and nonsingular matrices}
\begin{tarif}
The \textbf{rank} of an $m\times n$ matrix $A$ is the largest order of a
nonzero minor of $A$. It is denoted by $\operatorname{rank}A$,
$\operatorname{rg}A$ or $r(A)$.
\end{tarif}
Equivalently, the rank of a matrix is the maximal number of linearly
independent rows (or, which gives the same number, of linearly independent
columns). We shall use the following basic properties.
\begin{itemize}[leftmargin=1.3em]
\item $0 \leq \operatorname{rg}A \leq \min(m,n)$.
\item $\operatorname{rg}A = 0$ if and only if $A = O$ is the zero matrix.
\item $\operatorname{rg}A = \operatorname{rg}A^{T}$: the row rank equals the column rank.
\item Elementary transformations do not change the rank of a matrix.
\end{itemize}
\subsection{Computing the rank by means of minors}
\vspace{3pt}
\begin{tarif}
A matrix $A$ has rank $r$ if
\begin{enumerate}[label=\arabic*)]
\item $A$ has at least one \textbf{nonzero} minor of order $r$, and
\item every minor of order $r+1$ is zero (or there are no minors of that order).
\end{enumerate}
A nonzero minor of the highest order is called a \textbf{basis minor}.
\end{tarif}
\begin{misol}
Let us find the rank of
$A=\begin{pmatrix}1 & 2 & 3\\ 2 & 4 & 6\\ 0 & 1 & -1\end{pmatrix}$.
The only minor of order three is the determinant itself:
\[\kf{\det A = \begin{vmatrix}1&2&3\\2&4&6\\0&1&-1\end{vmatrix}}\]
\[\kf{= 1(-4-6) - 2(-2-0) + 3(2-0) = -10 + 4 + 6 = 0.}\]
Hence $\operatorname{rg}A < 3$. Among the minors of order two we find
\[\kf{M = \begin{vmatrix}1 & 2\\0 & 1\end{vmatrix} = 1\cdot 1 - 2\cdot 0 = 1 \neq 0,}\]
so $\operatorname{rg}A = 2$.
\end{misol}
\begin{izoh}
Since $A$ in this example is a $3\times 3$ matrix, its only minor of order
three is $\det A$, and the equality $\det A=0$ already shows that
$\operatorname{rg}A<3$.

To prove that $\operatorname{rg}A\ge 2$, a \emph{single} nonzero minor of
order two is enough; the remaining minors of order two need not be computed,
because they cannot change the rank.

If the first minor of order two we try turns out to be zero, we try another
one. If all minors of order two vanish and the matrix is not the zero matrix,
then $\operatorname{rg}A=1$.

For larger matrices checking every minor quickly becomes impractical (a
$4\times 4$ matrix already has $16$ minors of order three). In such cases one
uses Gaussian elimination or the method of bordering minors.
\end{izoh}
\begin{misol}
Find the rank of $B=\begin{pmatrix}2 & -1 & 3\\ 4 & -2 & 6\end{pmatrix}$.
Since $B$ is a $2\times 3$ matrix, $\operatorname{rg}B \leq 2$. We check all
minors of order two:
\[\kf{\begin{vmatrix}2&-1\\4&-2\end{vmatrix} = -4+4=0, \quad
\begin{vmatrix}2&3\\4&6\end{vmatrix} = 12-12=0,}\]
\[\kf{\begin{vmatrix}-1&3\\-2&6\end{vmatrix} = -6+6=0.}\]
All of them vanish, while the minor of order one $a_{11}=2$ is nonzero.
Therefore $\operatorname{rg}B = 1$. This is no surprise: the second row is
twice the first, $S_2 = 2S_1$ (here and below $S_i$ denotes the $i$th row).
\end{misol}
\bigskip
% =========================================================
\subsection{Computing the rank by elementary transformations (Gaussian elimination)}
\vspace{3pt}
\begin{tarif}
Using elementary row operations, we reduce the matrix to \textbf{row echelon
form}. The rank of the matrix equals the \textbf{number of nonzero rows} of
the resulting echelon matrix.
\end{tarif}
The \textbf{elementary row operations} are:
\begin{enumerate}[label=\arabic*)]
\item interchanging two rows, $S_i \leftrightarrow S_j$;
\item multiplying a row by a nonzero number, $S_i \to k\cdot S_i$, $k\neq 0$;
\item adding to a row a multiple of another row, $S_i \to S_i + k\cdot S_j$.
\end{enumerate}
A matrix is in \textbf{row echelon form} if the first nonzero entry of each
row (its \emph{leading entry}, or pivot) lies strictly to the right of the
leading entry of the row above it, and all zero rows stand at the bottom.
\vspace{3pt}
\begin{misol}
Find the rank of
$A=\begin{pmatrix}1 & 2 & -1 & 3\\ 2 & 5 & 0 & -1\\ -1 & 0 & 3 & 2\\ 3 & 1 & -2 & 4\end{pmatrix}$.
\begin{yechish}
We apply Gaussian elimination. Clearing the first column with the first row
and then the second column with the second row, we get
$$
\begin{pmatrix}1&2&-1&3\\2&5&0&-1\\-1&0&3&2\\3&1&-2&4\end{pmatrix}
\xrightarrow[\substack{S_3+S_1 \\ S_4-3S_1}]{S_2-2S_1}
\begin{pmatrix}1&2&-1&3\\0&1&2&-7\\0&2&2&5\\0&-5&1&-5\end{pmatrix}
$$
\[\kf{\xrightarrow[S_4+5S_2]{S_3-2S_2}
\begin{pmatrix}1&2&-1&3\\0&1&2&-7\\0&0&-2&19\\0&0&11&-40\end{pmatrix}
\xrightarrow{S_4+\frac{11}{2}S_3}
\begin{pmatrix}1&2&-1&3\\0&1&2&-7\\0&0&-2&19\\0&0&0&\frac{129}{2}\end{pmatrix}.}\]
The echelon form has four nonzero rows, so $\operatorname{rg}A = 4$
(indeed, $\det A=-129\neq 0$).
\end{yechish}
\end{misol}
\begin{misol}
Find the rank of
$B=\begin{pmatrix}1&-1&2&0\\2&-2&4&0\\3&1&-1&2\\1&3&-5&2\end{pmatrix}$.
\begin{yechish}
$$
\begin{pmatrix}1&-1&2&0\\2&-2&4&0\\3&1&-1&2\\1&3&-5&2\end{pmatrix}
\xrightarrow[\substack{S_3-3S_1 \\ S_4-S_1}]{S_2-2S_1}
\begin{pmatrix}1&-1&2&0\\0&0&0&0\\0&4&-7&2\\0&4&-7&2\end{pmatrix}
\xrightarrow{S_4-S_3}
$$
\[\kf{\begin{pmatrix}1&-1&2&0\\0&0&0&0\\0&4&-7&2\\0&0&0&0\end{pmatrix}}\]
\[\kf{\xrightarrow{S_2\leftrightarrow S_3}
\begin{pmatrix}1&-1&2&0\\0&4&-7&2\\0&0&0&0\\0&0&0&0\end{pmatrix}.}\]
Only two rows remain nonzero, hence $\operatorname{rg}B = 2$.
\end{yechish}
\end{misol}
\begin{misol}
Find the rank of $C=\begin{pmatrix}2&1&-3&5\\4&2&-6&10\\1&0&1&-2\end{pmatrix}$.
\begin{yechish}
It is convenient to bring a row with leading entry $1$ to the top:
\[\kf{\begin{pmatrix}2&1&-3&5\\4&2&-6&10\\1&0&1&-2\end{pmatrix}
\xrightarrow{S_1\leftrightarrow S_3}
\begin{pmatrix}1&0&1&-2\\4&2&-6&10\\2&1&-3&5\end{pmatrix}}\]
\[\kf{\xrightarrow[S_3-2S_1]{S_2-4S_1}
\begin{pmatrix}1&0&1&-2\\0&2&-10&18\\0&1&-5&9\end{pmatrix}
\xrightarrow{S_2-2S_3}
\begin{pmatrix}1&0&1&-2\\0&0&0&0\\0&1&-5&9\end{pmatrix}
\xrightarrow{S_2\leftrightarrow S_3}}\]
\[\kf{\begin{pmatrix}1&0&1&-2\\0&1&-5&9\\0&0&0&0\end{pmatrix}.}\]
There are two nonzero rows, so $\operatorname{rg}C = 2$.
\end{yechish}
\end{misol}
\textbf{Singular and nonsingular matrices.}
\vspace{3pt}
\begin{tarif}
A square matrix $A$ of order $n$ is called
\begin{itemize}[leftmargin=1.3em]
\item \textbf{nonsingular} (nondegenerate) if $\det A \neq 0$, that is,
$\operatorname{rg}A = n$;
\item \textbf{singular} (degenerate) if $\det A = 0$, that is,
$\operatorname{rg}A < n$.
\end{itemize}
\end{tarif}
A nonsingular matrix $A$ has the following properties:
\begin{enumerate}[label=\arabic*)]
\item it has an inverse $A^{-1}$, and the inverse is unique;
\item its rows (and its columns) are linearly independent;
\item the system $Ax=b$ has a unique solution for every right-hand side $b$;
\item $\operatorname{rg}A = n$ (full rank).
\end{enumerate}
For a singular matrix, on the contrary:
\begin{enumerate}[label=\arabic*)]
\item there is no inverse matrix;
\item its rows (columns) are linearly dependent;
\item the system $Ax=b$ is either inconsistent or has infinitely many solutions;
\item $\operatorname{rg}A < n$.
\end{enumerate}
\vspace{3pt}
\begin{misol}
For $A=\begin{pmatrix}1&2\\3&4\end{pmatrix}$ we have
$\det A = 1\cdot 4 - 2\cdot 3 = -2 \neq 0$, so $A$ is \textbf{nonsingular}
and $\operatorname{rg}A = 2$.
For $B=\begin{pmatrix}2&6\\1&3\end{pmatrix}$ we have
$\det B = 2\cdot 3 - 6\cdot 1 = 0$, so $B$ is \textbf{singular} and
$\operatorname{rg}B = 1$: the rows are proportional ($S_1=2S_2$), and so are
the columns (the second column is three times the first).
\end{misol}
\begin{misol}
Consider $C=\begin{pmatrix}1&0&2\\-1&3&1\\2&-3&1\end{pmatrix}$. Expanding
along the first row,
\[\kf{\det C = 1(3+3) - 0 + 2(3-6) = 6-6 = 0.}\]
Thus $C$ is \textbf{singular} and $\operatorname{rg}C < 3$. The minor of
order two in the upper left corner is
\[\kf{\begin{vmatrix}1&0\\-1&3\end{vmatrix}=3\neq 0,}\]
hence $\operatorname{rg}C = 2$.
\end{misol}
\textbf{Rank and systems of linear equations.}
\vspace{3pt}
\begin{teorema}[Kronecker–Capelli]
A system of linear equations $Ax=b$ is consistent if and only if
\[\kf{\operatorname{rg}A = \operatorname{rg}(A\,|\,b),}\]
that is, the rank of the coefficient matrix $A$ equals the rank of the
augmented matrix $(A\,|\,b)$. Moreover, if $n$ is the number of unknowns,
then
\begin{itemize}[leftmargin=1.3em]
\item for $\operatorname{rg}A = \operatorname{rg}(A\,|\,b) = n$ the system has
a \textbf{unique solution};
\item for $\operatorname{rg}A = \operatorname{rg}(A\,|\,b) = r < n$ it has
\textbf{infinitely many solutions}, depending on $n-r$ free unknowns;
\item for $\operatorname{rg}A \neq \operatorname{rg}(A\,|\,b)$ it has
\textbf{no solutions} (the system is inconsistent).
\end{itemize}
\end{teorema}
In the literature this result is also known as the Rouché–Capelli theorem.
\begin{misol}
Is the system
\[\kf{\begin{cases}x_1+2x_2-x_3=3\\2x_1+4x_2-2x_3=5\end{cases}}\]
consistent? Its coefficient matrix and augmented matrix are
\[\kf{A=\begin{pmatrix}1&2&-1\\2&4&-2\end{pmatrix}, \quad
(A\,|\,b)=\begin{pmatrix}1&2&-1&3\\2&4&-2&5\end{pmatrix}.}\]
The operation $S_2 - 2S_1$ gives
$(A\,|\,b) \to \begin{pmatrix}1&2&-1&3\\0&0&0&-1\end{pmatrix}$.
Thus $\operatorname{rg}A = 1$, whereas $\operatorname{rg}(A\,|\,b)=2$ because
the entry $-1$ survives in the last column. The ranks differ, so the system
\textbf{has no solutions}.
\end{misol}
\begin{misol}
Apply the Kronecker–Capelli theorem to the system
$\begin{cases}x_1+2x_2-x_3=3,\\2x_1+4x_2-2x_3=6,\\x_1+x_2+x_3=1.\end{cases}$
\begin{yechish}
Reducing the augmented matrix, we obtain
$(A\,|\,b)=\begin{pmatrix}1&2&-1&3\\2&4&-2&6\\1&1&1&1\end{pmatrix}
\xrightarrow[S_3-S_1]{S_2-2S_1}
\begin{pmatrix}1&2&-1&3\\0&0&0&0\\0&-1&2&-2\end{pmatrix}$.
Hence $\operatorname{rg}A=\operatorname{rg}(A\,|\,b)=2$. The ranks are
equal, but $r=2<n=3$, so the system has infinitely many solutions with
$n-r=1$ free unknown.
\end{yechish}
\end{misol}
\bigskip
% =========================================================
\subsection{Rank and singularity in information security}
\textbf{Nonsingular keys in cryptography.}
In the Hill cipher, as in other matrix-based ciphers, the key matrix must be
\textbf{nonsingular}. If the key $K$ were singular ($\det K = 0$,
$\operatorname{rg}K < n$), the inverse $K^{-1}$ would not exist and
decryption would be impossible. Worse, different plaintext blocks could be
mapped to the same ciphertext, which defeats the purpose of encryption. For
this reason the first thing one checks when choosing a key is the condition
$\det K \neq 0$ (for the Hill cipher, invertibility of $\det K$ modulo $26$).

\textbf{Rank in coding theory.}
In error-correcting codes (Hamming codes, Reed–Solomon codes and others) the
rank of the generator matrix $G$ determines how much information the code
carries. A $k\times n$ generator matrix must have rank $k$, that is,
linearly independent rows; if $\operatorname{rg}G < k$, two different
messages receive the same codeword and cannot be told apart. Likewise, the
parity-check matrix $H$ must have rank $n-k$.

\textbf{Rank analysis of networks.}
The rank of the Laplacian matrix of a network graph reflects its
\textbf{connectivity}. For a connected network with $n$ nodes the Laplacian
has rank $n-1$; a smaller rank means that the network splits into several
disconnected components, which is an important warning sign in a security
audit. The rank of the Laplacian and its null space also help to locate weak
points of the network.

\textbf{Lattice-based cryptography.}
In modern post-quantum schemes such as LWE (Learning With Errors), rank and
nonsingularity are again central. The public key contains a large random
matrix, and its having full rank is one of the requirements for the security
of the scheme. An attacker trying to recover the secret key exploits linear
dependences among the columns of this matrix, so the closer the matrix is to
full rank, the harder the attack.

\subsection{Computing with Python, Maple and Mathcad}
Let us check the rank of the matrix $B$ from the second example on
Gaussian elimination and the consistency test of the system
$x_1+2x_2-x_3=3$, $2x_1+4x_2-2x_3=5$. In Python we use SymPy, which works in
exact rational arithmetic; the method \texttt{rref} also returns the reduced
row echelon form.
\begin{Verbatim}[fontsize=\footnotesize,frame=single,label={\small Python}]
from sympy import Matrix
B = Matrix([[1, -1, 2, 0], [2, -2, 4, 0],
            [3, 1, -1, 2], [1, 3, -5, 2]])
print("rank B =", B.rank())
R, pivots = B.rref()
for row in R.tolist():
    print(row)
A = Matrix([[1, 2, -1], [2, 4, -2]])
b = Matrix([3, 5])
print(A.rank(), A.row_join(b).rank())
\end{Verbatim}
\begin{Verbatim}[fontsize=\footnotesize,frame=single,label={\small Output}]
rank B = 2
[1, 0, 1/4, 1/2]
[0, 1, -7/4, 1/2]
[0, 0, 0, 0]
[0, 0, 0, 0]
1 2
\end{Verbatim}
The last line shows $\operatorname{rg}A=1\neq\operatorname{rg}(A\,|\,b)=2$,
so the system is inconsistent. In Maple the same computation uses the
\texttt{LinearAlgebra} package:
\begin{Verbatim}[fontsize=\footnotesize,frame=single,label={\small Maple}]
with(LinearAlgebra):
B := Matrix([[1,-1,2,0], [2,-2,4,0],
             [3,1,-1,2], [1,3,-5,2]]):
Rank(B);                    # 2
ReducedRowEchelonForm(B);
A := Matrix([[1,2,-1], [2,4,-2]]):
b := Vector([3,5]):
Rank(A), Rank(<A|b>);       # 1, 2
\end{Verbatim}
In Mathcad Prime the matrix is entered through the matrix template; the
built-in functions \texttt{rank}, \texttt{rref} and \texttt{augment} do the
rest (Mathcad displays the fractions as decimals).
\begin{Verbatim}[fontsize=\footnotesize,frame=single,label={\small Mathcad}]
B := | 1  -1   2   0 |
     | 2  -2   4   0 |
     | 3   1  -1   2 |
     | 1   3  -5   2 |
rank(B) = 2
          | 1  0   0.25  0.5 |
rref(B) = | 0  1  -1.75  0.5 |
          | 0  0   0     0   |
          | 0  0   0     0   |
A := | 1  2  -1 |      b := | 3 |
     | 2  4  -2 |           | 5 |
rank(A) = 1     rank(augment(A, b)) = 2
\end{Verbatim}

\subsection{Exercises on the rank of a matrix}
\bigskip
\emph{1. Find the rank of each matrix by means of minors.}
\bigskip
\begin{multicols}{3}\raggedright
\begin{enumerate}[label=\arabic*.,  leftmargin=*]
\item $\begin{pmatrix}1&2&3\\2&4&6\\0&1&-1\end{pmatrix}$
\item $\begin{pmatrix}2&-1&3\\4&-2&6\\1&0&-1\end{pmatrix}$
\item $\begin{pmatrix}3&1&-2\\6&2&-4\\-3&-1&2\end{pmatrix}$
\item $\begin{pmatrix}1&0&2\\-1&3&1\\2&-3&1\end{pmatrix}$
\item $\begin{pmatrix}0&1&-3\\2&-1&5\\4&-2&10\end{pmatrix}$
\item $\begin{pmatrix}5&-1&2\\3&4&0\\-2&9&-2\end{pmatrix}$
\item $\begin{pmatrix}1&-2&3\\-2&4&-6\\3&-6&9\end{pmatrix}$
\item $\begin{pmatrix}2&1&0\\0&3&-1\\4&5&-1\end{pmatrix}$
\item $\begin{pmatrix}-1&3&2\\2&-6&-4\\0&0&0\end{pmatrix}$
\item $\begin{pmatrix}4&-2&1\\0&3&-5\\8&-4&2\end{pmatrix}$
\item $\begin{pmatrix}1&1&1\\2&2&2\\3&3&3\end{pmatrix}$
\item $\begin{pmatrix}3&0&-1\\0&2&4\\-6&0&2\end{pmatrix}$
\item $\begin{pmatrix}1&2&-1\\3&1&0\\5&5&-2\end{pmatrix}$
\item $\begin{pmatrix}-2&4&1\\1&-2&3\\0&0&7\end{pmatrix}$
\item $\begin{pmatrix}0&0&5\\0&3&-1\\2&0&0\end{pmatrix}$
\item $\begin{pmatrix}1&3&5\\2&6&10\\-1&-3&-5\end{pmatrix}$
\item $\begin{pmatrix}4&-1&2\\-8&2&-4\\2&3&-1\end{pmatrix}$
\item $\begin{pmatrix}0&1&2\\1&0&-3\\2&-1&0\end{pmatrix}$
\item $\begin{pmatrix}3&6&-3\\1&2&-1\\-2&-4&2\end{pmatrix}$
\item $\begin{pmatrix}-1&0&4\\2&-3&1\\1&-3&5\end{pmatrix}$
\item $\begin{pmatrix}5&10&-5\\1&2&-1\\0&0&0\end{pmatrix}$
\item $\begin{pmatrix}2&-1&0\\-4&2&0\\1&3&-5\end{pmatrix}$
\item $\begin{pmatrix}1&0&0\\0&1&0\\0&0&0\end{pmatrix}$
\item $\begin{pmatrix}3&-2&1\\0&0&0\\-6&4&-2\end{pmatrix}$
\item $\begin{pmatrix}1&-1&2\\2&0&1\\3&-1&3\end{pmatrix}$
\item $\begin{pmatrix}-3&1&5\\6&-2&-10\\0&4&-1\end{pmatrix}$
\item $\begin{pmatrix}2&4&-2\\-1&-2&1\\3&6&-3\end{pmatrix}$
\item $\begin{pmatrix}0&3&1\\-2&0&4\\2&-3&3\end{pmatrix}$
\item $\begin{pmatrix}1&-3&0\\-2&6&0\\4&-12&0\end{pmatrix}$
\item $\begin{pmatrix}7&0&-1\\0&5&2\\-14&0&2\end{pmatrix}$
\end{enumerate}
\end{multicols}
\emph{2. Find the rank of each matrix by Gaussian elimination (elementary row operations).}
\bigskip
\begin{multicols}{2}\raggedright
\begin{enumerate}[label=\arabic*.,  leftmargin=*]
\item $\begin{pmatrix}1&-1&2&0\\2&-2&4&0\\3&1&-1&2\end{pmatrix}$
\item $\begin{pmatrix}2&1&-3&5\\4&2&-6&10\\1&0&1&-2\end{pmatrix}$
\item $\begin{pmatrix}1&2&3&4\\2&4&6&8\\3&6&9&12\end{pmatrix}$
\item $\begin{pmatrix}1&0&-1&2\\3&1&0&-1\\2&1&1&-3\end{pmatrix}$
\item $\begin{pmatrix}0&1&2&-3\\1&0&-1&4\\2&1&0&5\end{pmatrix}$
\item $\begin{pmatrix}1&3&-2\\4&12&-8\\-2&-6&4\end{pmatrix}$
\item $\begin{pmatrix}3&-1&2&0\\6&-2&4&0\\1&1&-1&3\\4&0&1&3\end{pmatrix}$
\item $\begin{pmatrix}1&2&-1\\0&3&4\\-1&1&5\\2&7&2\end{pmatrix}$
\item $\begin{pmatrix}2&0&1&-3\\4&0&2&-6\\1&-1&0&2\\3&-1&1&-1\end{pmatrix}$
\item $\begin{pmatrix}1&-2&3\\-2&4&-6\\5&-10&15\end{pmatrix}$
\item $\begin{pmatrix}0&0&1&2\\1&0&0&-1\\0&1&0&3\\1&1&1&2\end{pmatrix}$
\item $\begin{pmatrix}1&1&-1&0\\2&-1&0&3\\3&0&-1&3\end{pmatrix}$
\item $\begin{pmatrix}5&-2&1\\-10&4&-2\\15&-6&3\end{pmatrix}$
\item $\begin{pmatrix}1&0&2&1\\0&1&-1&3\\1&1&1&4\\2&1&3&5\end{pmatrix}$
\item $\begin{pmatrix}-1&3&0&2\\2&-6&0&-4\\0&0&1&5\end{pmatrix}$
\item $\begin{pmatrix}1&2&3\\4&5&6\\7&8&9\end{pmatrix}$
\item $\begin{pmatrix}2&-1&0&4\\0&3&-2&1\\4&-2&0&8\\2&2&-2&5\end{pmatrix}$
\item $\begin{pmatrix}1&-1&2&-3\\2&-2&4&-6\\-1&1&-2&3\end{pmatrix}$
\item $\begin{pmatrix}0&2&-1\\3&0&4\\6&-4&9\\3&2&3\end{pmatrix}$
\item $\begin{pmatrix}1&0&0&1\\0&1&0&1\\0&0&1&1\\1&1&1&0\end{pmatrix}$
\item $\begin{pmatrix}3&1&-2&0\\0&0&0&0\\6&2&-4&0\end{pmatrix}$
\item $\begin{pmatrix}1&-3&2\\-2&6&-4\\4&-12&8\\0&0&0\end{pmatrix}$
\item $\begin{pmatrix}2&1&0&-1\\0&-2&3&4\\2&-1&3&3\end{pmatrix}$
\item $\begin{pmatrix}1&1&1\\1&2&3\\1&3&6\end{pmatrix}$
\item $\begin{pmatrix}0&1&-2&3\\0&2&-4&6\\1&0&1&-1\\2&-1&4&-5\end{pmatrix}$
\item $\begin{pmatrix}4&-2&6\\2&-1&3\\-6&3&-9\end{pmatrix}$
\item $\begin{pmatrix}1&0&-1\\0&1&2\\1&1&1\\2&1&0\end{pmatrix}$
\item $\begin{pmatrix}3&-1&2&1\\1&0&-1&2\\5&-1&1&3\\4&-1&1&3\end{pmatrix}$
\item $\begin{pmatrix}1&2&0\\-1&0&2\\0&2&2\\2&4&0\end{pmatrix}$
\item $\begin{pmatrix}0&0&0\\1&-2&3\\-2&4&-6\\3&-6&9\end{pmatrix}$
\end{enumerate}
\end{multicols}
\emph{3. Determine whether each matrix is singular or nonsingular (compute $\det A$).}
\bigskip
\begin{multicols}{3}\raggedright
\begin{enumerate}[label=\arabic*.,  leftmargin=*]
\item $\begin{pmatrix}1&2\\3&4\end{pmatrix}$
\item $\begin{pmatrix}2&6\\1&3\end{pmatrix}$
\item $\begin{pmatrix}5&-1\\3&2\end{pmatrix}$
\item $\begin{pmatrix}-3&9\\1&-3\end{pmatrix}$
\item $\begin{pmatrix}4&0\\0&-2\end{pmatrix}$
\item $\begin{pmatrix}1&0&2\\-1&3&1\\2&-3&1\end{pmatrix}$
\item $\begin{pmatrix}3&-1&2\\0&5&-4\\6&-2&4\end{pmatrix}$
\item $\begin{pmatrix}2&1&-3\\4&2&-6\\1&0&1\end{pmatrix}$
\item $\begin{pmatrix}1&-2&3\\0&4&-1\\2&0&5\end{pmatrix}$
\item $\begin{pmatrix}-1&3&0\\2&-1&4\\5&5&4\end{pmatrix}$
\item $\begin{pmatrix}0&1\\-1&0\end{pmatrix}$
\item $\begin{pmatrix}7&14\\2&4\end{pmatrix}$
\item $\begin{pmatrix}-2&0&5\\1&3&-4\\0&6&-3\end{pmatrix}$
\item $\begin{pmatrix}1&1&1\\1&2&3\\1&3&6\end{pmatrix}$
\item $\begin{pmatrix}3&6&9\\1&2&3\\-2&-4&-6\end{pmatrix}$
\item $\begin{pmatrix}8&-4\\-6&3\end{pmatrix}$
\item $\begin{pmatrix}5&0&0\\0&-3&0\\0&0&2\end{pmatrix}$
\item $\begin{pmatrix}1&2&3\\4&5&6\\5&7&9\end{pmatrix}$
\item $\begin{pmatrix}-1&0\\0&-1\end{pmatrix}$
\item $\begin{pmatrix}2&-3&1\\4&-6&2\\-1&0&3\end{pmatrix}$
\item $\begin{pmatrix}10&5\\-2&-1\end{pmatrix}$
\item $\begin{pmatrix}0&0&1\\0&2&0\\3&0&0\end{pmatrix}$
\item $\begin{pmatrix}1&-1&2\\3&0&-4\\4&-1&-2\end{pmatrix}$
\item $\begin{pmatrix}6&2\\9&3\end{pmatrix}$
\item $\begin{pmatrix}2&0&-1\\0&3&0\\4&0&-2\end{pmatrix}$
\item $\begin{pmatrix}-5&1&0\\3&-2&4\\-7&3&-4\end{pmatrix}$
\item $\begin{pmatrix}1&0\\0&0\end{pmatrix}$
\item $\begin{pmatrix}3&1&-2\\-6&-2&4\\9&3&-6\end{pmatrix}$
\item $\begin{pmatrix}4&-3\\-8&6\end{pmatrix}$
\item $\begin{pmatrix}1&2&0\\0&1&3\\-1&0&6\end{pmatrix}$
\end{enumerate}
\end{multicols}

\emph{4. Using the Kronecker–Capelli theorem, find the ranks of the coefficient matrix and of the augmented matrix, and decide whether the system is consistent.}
\bigskip
\begin{multicols}{2}\raggedright
\small
\begin{enumerate}[label=\arabic*.,  leftmargin=*]
\item $\begin{cases}x_1+2x_2-x_3=3\\2x_1+4x_2-2x_3=6\\x_1+x_2+x_3=1\end{cases}$
\item $\begin{cases}x_1-x_2+2x_3=1\\2x_1-2x_2+4x_3=3\end{cases}$
\item $\begin{cases}x_1+x_2+x_3=1\\2x_1+3x_2-x_3=4\\3x_1+4x_2=5\end{cases}$
\item $\begin{cases}2x_1-x_2=3\\-4x_1+2x_2=-6\end{cases}$
\item $\begin{cases}x_1+2x_2+3x_3=1\\4x_1+5x_2+6x_3=2\\7x_1+8x_2+9x_3=3\end{cases}$
\item $\begin{cases}x_1-3x_2=5\\2x_1-6x_2=7\end{cases}$
\item $\begin{cases}x_1+x_2-x_3+x_4=2\\2x_1+2x_2-2x_3+2x_4=4\\x_1-x_2+x_3=1\end{cases}$
\item $\begin{cases}3x_1+x_2-2x_3=1\\6x_1+2x_2-4x_3=5\end{cases}$
\item $\begin{cases}x_1+2x_2=5\\3x_1-x_2=1\\5x_1+3x_2=11\end{cases}$
\item $\begin{cases}x_1-x_2+x_3=0\\2x_1+x_2-x_3=0\\x_1+2x_2-2x_3=0\end{cases}$
\item $\begin{cases}x_1+x_2=3\\2x_1+2x_2=6\\3x_1+3x_2=9\end{cases}$
\item $\begin{cases}2x_1-x_2+3x_3=4\\4x_1-2x_2+6x_3=7\end{cases}$
\item $\begin{cases}x_1+3x_2-2x_3=0\\2x_1+6x_2-4x_3=0\end{cases}$
\item $\begin{cases}x_1-x_2+x_3-x_4=1\\2x_1+x_2-x_3+2x_4=3\\x_1+2x_2-2x_3+3x_4=2\end{cases}$
\item $\begin{cases}x_1+2x_2+x_3=4\\2x_1+4x_2+2x_3=8\\-x_1-2x_2-x_3=-4\end{cases}$
\item $\begin{cases}5x_1-x_2=10\\-10x_1+2x_2=-15\end{cases}$
\item $\begin{cases}x_1+x_2+x_3=6\\x_1-x_2+x_3=2\\x_1+x_2-x_3=4\end{cases}$
\item $\begin{cases}x_1+2x_2-3x_3=1\\2x_1+4x_2-6x_3=2\\-x_1-2x_2+3x_3=-1\end{cases}$
\item $\begin{cases}x_1-2x_2=0\\3x_1+x_2=0\\2x_1+3x_2=0\end{cases}$
\item $\begin{cases}x_1+x_2+x_3+x_4=0\\2x_1-x_2+3x_3=5\\x_1+2x_2-x_3+4x_4=-5\end{cases}$
\item $\begin{cases}2x_1+3x_2=7\\4x_1+6x_2=14\end{cases}$
\item $\begin{cases}x_1-x_2+2x_3=3\\3x_1-3x_2+6x_3=9\\-2x_1+2x_2-4x_3=-6\end{cases}$
\item $\begin{cases}x_1+2x_2-x_3=1\\3x_1+x_2+2x_3=4\\5x_1+5x_2=6\end{cases}$
\item $\begin{cases}x_1-x_2=2\\2x_1+3x_2=1\\4x_1+x_2=5\end{cases}$
\item $\begin{cases}x_1+x_2+x_3=1\\2x_1+2x_2+2x_3=2\\3x_1+3x_2+3x_3=3\end{cases}$
\item $\begin{cases}3x_1-2x_2+x_3=4\\6x_1-4x_2+2x_3=1\end{cases}$
\item $\begin{cases}x_1+x_2=1\\x_1-x_2=3\\2x_1+x_2=4\end{cases}$
\item $\begin{cases}x_1+2x_2+3x_3=0\\2x_1+4x_2+6x_3=0\end{cases}$
\item $\begin{cases}x_1-x_2+x_3=2\\2x_1-2x_2+2x_3=4\\3x_1-3x_2+3x_3=7\end{cases}$
\item $\begin{cases}x_1+x_2-x_3=3\\2x_1-x_2+x_3=1\\3x_1+x_2-x_3=5\end{cases}$
\end{enumerate}
\end{multicols}
